\chapter[1919 --- Bordet]{1919: Jules Bordet}
\label{ch:1919_julesbordet}

\section*{Main Contribution}

\textit{Discovered complement, a system of blood proteins that enhance the ability of antibodies and phagocytic cells to clear pathogens---a key component of innate immunity.}

\section{Official Citation}

for his discoveries relating to immunity

\section{Background and Historical Context}

Jules Bordet was born on June 13, 1870, in Soignies, a small town in the province of Hainaut, Belgium. He was the son of a schoolmaster and grew up in a modest but intellectually stimulating household that valued education and scientific inquiry. From an early age, Bordet displayed a keen curiosity about the natural world, and after completing his secondary education, he enrolled at the Free University of Brussels to study medicine. It was there that he encountered the emerging discipline of bacteriology, which was being revolutionized by the discoveries of Louis Pasteur and Robert Koch. In 1892, shortly after completing his medical degree, Bordet traveled to Paris to work in the laboratory of \'{E}lie Metchnikoff at the newly established Pasteur Institute---an experience that would shape the trajectory of his entire career.

Metchnikoff, who would himself share the Nobel Prize in 1908, was a charismatic and innovative scientist who had proposed the theory of phagocytosis in the 1880s. This theory held that certain white blood cells---later named macrophages and neutrophils---were capable of engulfing and destroying invading microorganisms, and that this cellular mechanism constituted the primary defense of the body against infection. Metchnikoff's ideas were controversial and were vigorously opposed by the school of Paul Ehrlich, the German physician and scientist who argued that immunity was mediated by soluble substances in the blood---what we now call antibodies---and not by cells. This great debate between the ``cellularists'' and the ``humoralists'' dominated immunology in the final years of the nineteenth century, and it was in this intellectually charged atmosphere that Bordet made his major discoveries.

During his time at the Pasteur Institute, Bordet conducted experiments that would prove crucial to resolving the debate between Metchnikoff and Ehrlich. He demonstrated that the destruction of bacteria by immune serum required the cooperation of both cellular and humoral components---that is, both antibodies and phagocytic cells played essential and complementary roles in the immune response. This synthesis was characteristic of Bordet's approach to science: rather than taking sides in ideological disputes, he sought to understand the actual mechanisms at work, even when those mechanisms were more complex and nuanced than either side had imagined.

After returning to Belgium in 1895, Bordet was appointed director of the newly established Anti-Rabies Institute in Brussels, a position that gave him both the laboratory space and the clinical material needed to pursue his research program. Over the next two decades, Bordet conducted a series of meticulous experiments on the mechanisms of immunity, focusing particularly on the phenomenon of bacterial lysis---the dissolution of bacterial cells by immune serum. His work in this period would fundamentally reshape our understanding of the immune system and would lead to the development of diagnostic tools that saved countless lives.

The scientific community of the early 1900s was deeply divided over the mechanisms by which the body destroyed invading microorganisms. The humoralists, led by Ehrlich, argued that soluble substances in the blood---antibodies---were the primary mediators of immunity, and that these antibodies could be produced in large quantities in response to immunization. The cellularists, following Metchnikoff, maintained that phagocytic cells were the main line of defense, and that antibodies played only a supporting role in enhancing phagocytosis. Bordet's genius lay in his ability to transcend this dichotomy and to demonstrate that both humoral and cellular mechanisms worked together in a coordinated fashion to eliminate pathogens---a synthesis that was not fully appreciated until decades later but that now forms the foundation of modern immunology.

\section{The Discovery and Contribution}

Bordet's most celebrated contribution to immunology was his elucidation of the role of complement in immune defense, and the development of the complement fixation technique. To understand the significance of this work, it is necessary to appreciate the state of knowledge at the time. Scientists knew that serum from immunized animals could kill bacteria, but the mechanism by which this occurred was poorly understood. Bordet conducted a series of elegant experiments that dissected this process into its component parts, revealing for the first time the multi-step architecture of the humoral immune response.

Working with the bacterium \emph{Vibrio cholerae} and other organisms, Bordet showed that when bacteria were mixed with serum from an immunized animal, the bacteria were first ``sensitized''---that is, they were coated with a specific substance (antibody) that rendered them susceptible to lysis. However, this sensitization alone was insufficient to kill the bacteria. A second substance, present in fresh serum but destroyed by heating to 56\,$^{\circ}$C for 30 minutes, was also required. Bordet called this second substance ``alexin,'' a term that was later replaced by ``complement'' (introduced by Paul Ehrlich, who used the term to denote a substance that ``completed'' the lytic reaction). The complement system, as we now know it, consists of a cascade of over 30 proteins---including components C1 through C9, factor B, factor D, properdin, and many regulatory proteins---that work together to lyse bacteria, opsonize pathogens for phagocytosis, and recruit inflammatory cells to the site of infection.

Bordet's experiments were remarkably precise and methodologically innovative. He demonstrated that complement was not specific---it could lyse any sensitized bacterium, regardless of the antibody specificity. This was in contrast to antibodies, which were highly specific for particular bacterial antigens. This distinction between the specific and non-specific components of humoral immunity was a conceptual breakthrough that clarified the architecture of the immune response. Bordet also showed that complement was consumed during the lytic reaction---a finding that had significant implications for the development of diagnostic tests, because it meant that the amount of complement consumed could be used as a measure of the amount of antigen-antibody reaction that had occurred.

Building on these fundamental discoveries, Bordet developed the complement fixation test (also known as the Bordet-Gengou test, after his collaborator Octave Gengou). The principle of the test was elegant in its simplicity: if a patient's serum contained antibodies against a particular pathogen, those antibodies would fix complement when mixed with the corresponding antigen. After incubation, the mixture would be tested for residual complement activity by adding sensitized red blood cells (red cells that had been coated with anti-red cell antibody). If complement had been fixed by the antigen-antibody reaction, the red blood cells would not be lysed (because the complement had been consumed); if no antibodies were present, complement would remain free and would lyse the indicator red cells, producing a visible hemolysis.

The complement fixation test represented a quantum leap in diagnostic medicine. For the first time, physicians had a reliable method for detecting the presence of specific antibodies in a patient's blood, thereby diagnosing infectious diseases that had previously been difficult or impossible to identify during the early stages of infection. The test was quickly adapted for the diagnosis of syphilis (the Wassermann test, developed by August von Wassermann in 1906 but based on Bordet's complement fixation principles), tuberculosis, gonorrhea, and many other infectious diseases. It remained the gold standard for serological diagnosis for decades, and its fundamental principle---the detection of antigen-antibody reactions through a measurable indicator system---is still employed in modern immunoassays.

Beyond complement, Bordet made numerous other contributions to immunology and microbiology. He demonstrated that the agglutination of bacteria by immune serum was a specific reaction---the basis of the agglutination test, which became one of the most widely used methods for the serological identification of bacteria. He showed that the sensitization of bacteria by antibody was an irreversible process, and that the affinity between antibody and antigen was extremely high. He also contributed to the understanding of opsonins---serum factors that enhanced the phagocytosis of bacteria by white blood cells---bridging the gap between humoral and cellular immunity by showing that antibodies and complement proteins worked together to coat pathogen surfaces and facilitate their uptake by phagocytes.

In later years, Bordet turned his attention to the mechanisms of blood coagulation, demonstrating that the clotting of blood involved a series of enzymatic reactions triggered by the release of thromboplastin (now called tissue factor) from damaged tissues. He showed that blood coagulation required the sequential activation of a series of plasma clotting factors, and that the process could be initiated by either the intrinsic pathway (triggered by contact with exposed collagen) or the extrinsic pathway (triggered by tissue factor release). This work anticipated the modern understanding of the coagulation cascade, which is now recognized as one of the most complex and finely regulated enzymatic systems in the body.

Bordet was also a gifted teacher and science communicator. His textbook, \emph{Trait{\'e} de l'immunit{\'e} dans les maladies infectieuses} (Treatise on Immunity in Infectious Diseases), published in 1920, became the definitive reference in the field and was translated into multiple languages. The book was notable for its clarity and its integration of experimental findings with clinical applications, reflecting Bordet's conviction that basic research and medical practice were inseparable. Bordet was elected to the Belgian Senate in 1922 and served as Minister of Education, in which capacity he championed the cause of scientific research and medical education in Belgium.

The Nobel Committee awarded Bordet the Prize in Physiology or Medicine in 1919 ``in recognition of his work on immunity,'' although the actual ceremony was delayed until 1920 due to the aftermath of World War I. The award was richly deserved: Bordet's discoveries had not only advanced fundamental understanding of the immune system but had also provided practical tools that saved countless lives through improved diagnosis of infectious diseases.

\section{Impact on Modern Medicine}

The impact of Bordet's work on modern medicine cannot be overstated. The complement system that he discovered and characterized is now recognized as a central component of innate immunity, playing critical roles in pathogen defense, inflammation, and tissue homeostasis. The complement system operates through three main activation pathways---the classical pathway (triggered by antigen-antibody complexes), the alternative pathway (triggered by direct contact with pathogen surfaces), and the lectin pathway (triggered by mannose-binding lectin binding to microbial carbohydrates)---all of which converge on a central event: the cleavage of complement component C3 into C3a and C3b. This central event triggers a cascade of downstream reactions that result in the formation of the membrane attack complex (C5b-C9), which creates pores in the membranes of target cells, leading to their lysis and death.

Deficiencies in complement components are associated with increased susceptibility to infections (particularly with encapsulated bacteria such as \emph{Streptococcus pneumoniae} and \emph{Neisseria meningitidis}), autoimmune diseases such as systemic lupus erythematosus (SLE), and atypical hemolytic uremic syndrome (aHUS). The therapeutic manipulation of complement---both inhibition and activation---is an active area of drug development, with monoclonal antibodies targeting complement proteins already approved for the treatment of conditions such as paroxysmal nocturnal hemoglobinuria (PNH, treated with eculizumab, an anti-C5 monoclonal antibody) and age-related macular degeneration (treated with ranibizumab, which inhibits the complement regulator C5a).

The complement fixation test, while largely superseded by newer immunoassay technologies such as ELISA and chemiluminescent immunoassays, remains a foundational concept in serology. The enzyme-linked immunosorbent assay (ELISA), which revolutionized diagnostic medicine in the 1970s and 1980s and remains one of the most widely used diagnostic tests in the world, is a direct descendant of the complement fixation test in its fundamental logic: the detection of specific antibodies or antigens in patient samples through carefully designed immunological reactions that produce a measurable signal. Modern multiplex bead-based immunoassays (such as Luminex), lateral flow assays used in point-of-care testing (such as rapid antigen tests for infectious diseases), and high-throughput automated immunoassay platforms all trace their conceptual lineage to Bordet's pioneering work on complement fixation and antigen-antibody reactions.

Bordet's demonstration that immune defense involves the cooperation of multiple effector mechanisms---antibodies, complement, and phagocytes---established the paradigm of integrated immune responses that guides immunology to this day. The modern understanding of opsonization, in which antibodies (particularly IgG) and complement proteins (particularly C3b) coat pathogen surfaces to facilitate phagocytic uptake by macrophages and neutrophils bearing Fc receptors and complement receptors, is a direct extension of Bordet's observations. This tripartite model of humoral and cellular cooperation remains the framework within which we understand host defense against bacterial infections, and it informs the design of vaccines that aim to elicit both strong antibody responses and effective phagocytic killing.

Furthermore, Bordet's contributions to the understanding of blood coagulation remain relevant. His identification of the tissue factor pathway and the role of thromboplastin in initiating clot formation anticipated modern coagulation cascade models. The diagnostic assays used to evaluate bleeding disorders---including the prothrombin time (PT), the activated partial thromboplastin time (aPTT), and the thrombin time---all build upon the foundational principles that Bordet elucidated. The development of anticoagulant drugs, such as warfarin (which inhibits vitamin K-dependent clotting factors) and the direct oral anticoagulants (which inhibit thrombin or factor Xa), also relies on an understanding of the coagulation cascade that can be traced back to Bordet's early work.

\section{Legacy}

Jules Bordet's legacy extends far beyond the complement system and the complement fixation test. He exemplified the ideal of the physician-scientist---a clinician who used the tools of basic science to address pressing medical problems. His career at the Brussels school of medicine inspired generations of Belgian and international scientists, and the institute he founded (the Jules Bordet Institute, now a comprehensive cancer center in Brussels) remains one of Europe's leading centers for medical research and treatment. The Bordet Medal, awarded by the Belgian Royal Academy of Medicine, recognizes outstanding contributions to the biomedical sciences in his memory.

Bordet's intellectual honesty and methodological rigor set standards that remain exemplary. He was careful to distinguish between observation and interpretation, and he resisted the temptation to construct elaborate theoretical edifices on incomplete evidence. In an era when immunology was rife with competing theories and bitter personal rivalries---the disputes between Ehrlich and Metchnikoff were legendary---Bordet maintained a measured and objective approach that earned him the respect of colleagues worldwide. His approach to science, combining meticulous experimentation with clear thinking and honest reporting, remains a model for investigators in all fields.

The concepts that Bordet pioneered---complement activation, opsonization, immune complex formation, and serological diagnosis---are woven into the fabric of modern medicine. Every time a physician orders a blood test to detect antibodies (as in the diagnosis of HIV, hepatitis, or autoimmune diseases), every time a vaccine is designed to elicit complement-fixing antibodies (as in the design of conjugate vaccines against encapsulated bacteria), and every time a researcher studies the complement cascade to develop new therapeutics (as in the development of complement inhibitors for autoimmune diseases), they are building upon the foundations that Jules Bordet laid more than a century ago. His work reminds us that fundamental discoveries in basic science, pursued with patience and rigor, can transform the practice of medicine in ways that the discoverer could scarcely have imagined. Bordet died in Brussels on April 6, 1961, at the age of 90, leaving behind a scientific legacy that continues to illuminate the mechanisms of immunity and to guide the development of new diagnostic and therapeutic approaches to infectious and autoimmune diseases.


\section{Modern Developments}

Bordet's complement research has evolved into modern complement-targeted therapies. Eculizumab, a monoclonal antibody targeting complement component C5, treats paroxysmal nocturnal hemoglobinuria and atypical hemolytic uremic syndrome. Understanding of complement dysregulation has revealed its role in age-related macular degeneration, Alzheimer's disease, and COVID-19. Complement-based diagnostics, including C3 and C4 assays, remain essential for evaluating immune deficiencies and autoimmune diseases.
